\section{Архитектура модели}

Модель основана на архитектуре энкодер-декодер (рис.~\ref{fig:model_block_diagram}), выполняющей преобразование ${C_{\text{in}}}$-канального тензора длиной ${L_1}$ в ${C_{\text{out}}}$-канальный тензор той же длины:
\begin{equation*}
    \textbf{Y}: \, \mathbb{R}^{C_{\text{in}} \times L_1} \mapsto \mathbb{R}^{C_{\text{out}} \times L_1}.
\end{equation*}

\begin{figure}[H]
	\centering
	\includegraphics [width=1.00\linewidth] {parametersEstimation/model_block_diagram.drawio.png}

	\caption{Блок-схема архитектуры модели}
	\label{fig:model_block_diagram}
\end{figure}

Преобразование ${\textbf{Y}}$ выполняется в 3 основных этапа:
\begin{itemize}
    \item Кодирование в скрытое (латентное) пространство (CNN-энкодер):
        \begin{equation*}
            \textbf{E}: \, \mathbb{R}^{F_0 \times L_1} \mapsto \mathbb{R}^{F_3 \times L_3}.
        \end{equation*}
    \item Выделение наиболее информативных признаков в скрытом пространстве (\enquote{Bottleneck}):
        \begin{equation*}
            \textbf{T}: \, \mathbb{R}^{F_3 \times L_3} \mapsto \mathbb{R}^{F_4 \times L_4} \mapsto \mathbb{R}^{F_4 \times L_4}.
        \end{equation*}
    \item Декодирование из скрытого пространства (CNN-декодер):
        \begin{equation*}
            \textbf{D}: \, \mathbb{R}^{F_4 \times L_4} \mapsto \mathbb{R}^{F_1 \times L_1}.
        \end{equation*}
\end{itemize}

Операции свертки на входе и выходе модели служат для согласования числа каналов (линейная проекция):
\begin{equation*}
    \begin{aligned}
        \mathbb{R}^{C_\text{in} \times L_1} &\mapsto \mathbb{R}^{F_0 \times L_1}, \\
        \mathbb{R}^{F_0 \times L_1} &\mapsto \mathbb{R}^{C_\text{out} \times L_1}.
    \end{aligned}
\end{equation*}

Аналогичную функцию проекции выполняет свертка на входе алгоритма \enquote{Attention} в \enquote{Bottleneck}, но уже с нормализацией и активацией:
\begin{equation*}
    \mathbb{R}^{F_3 \times L_4} \mapsto \mathbb{R}^{F_4 \times L_4}.
\end{equation*}

Выходной тензор ${\mathbb{R}^{C_\text{out} \times L_1}}$ представляет собой логиты ${d_i}$, ${i \in \{1, 2, \dots, L_1\}}$, преобразуемые сначала в вероятности ${p_i}$, затем в бинарные маски ${m_i}$:
\begin{equation*}
    \begin{aligned}
        p_i &= \frac{1}{1 + e^{-d_i}},\\
        y_i &= \mathbb{I}\!\left( p_i > \tau \right),
    \end{aligned}
\end{equation*}
где ${\forall i \colon y_i \in \{0, 1\}}$,
\begin{equation*}
    \mathbb{I}(p_i \geq \tau) = 
        \begin{cases}
            1, & p_i \geq \tau, \\
            0, & p_i < \tau,
        \end{cases}
\end{equation*}
${\tau}$ -- пороговый уровень, ${\tau \in \left[0; 1 \right]}$.

Каждый канал маски соответствует типу детектируемых объектов.

Количество обучаемых параметров построенной модели составляет порядка \num{5.2}~M.

\subsection{Энкодер}
Энкодер преобразует входные данные в латентное (скрытое) пространство уменьшенной размерности. В созданной реализации энкодер состоит из каскада $N$ сверточных базовых блоков. Заложена \enquote{residual}-архитектура базового блока по типу \enquote{pre-activation} (рис.~\ref{fig:residual_conv_diagram}), обеспечивающая лучший градиентный поток для быстрой сходимости при обучении \cite{he2016identity}.

\begin{figure}[H]
	\centering
	\includegraphics [width=0.75\linewidth] {parametersEstimation/residual_conv_diagram.drawio.png}

	\caption{Блок-схема базового блока свертки в энкодере и декодере}
	\label{fig:residual_conv_diagram}
\end{figure}

Каждый $n$-й базовы блок, ${n \in \{1, \dots, N\}}$, выполняет последовательно две свертки с ядрами ${\mathbf{W}_n^{1\text{enc}} \in \mathbb{R}^{F_{n-1} \times F_n}}$, ${\mathbf{W}_n^{2\text{enc}} \in \mathbb{R}^{F_n \times F_n}}$, соответственно:

\begin{equation*}
    \begin{aligned}
        \mathbf{z}_n^{1\text{enc}} &= \mathbf{W}_n^{1\text{enc}} \ast \sigma\left(\mathcal{N}_n^{1\text{enc}}\left(\mathbf{e}_{n-1}\right)\right), \\
        \mathbf{z}_n^{2\text{enc}} &= \mathbf{W}_n^{2\text{enc}} \ast \sigma\left(\mathcal{N}_n^{2\text{enc}}\left(\mathbf{z}_n^{1\text{enc}}\right)\right), \\
        \mathbf{e}_{n} &= \mathbf{z}_n^{2\text{enc}} \oplus \mathbf{s}_n^\text{enc},
    \end{aligned}
\end{equation*}
где \enquote{${\ast}$} -- операция свертки, \enquote{${\oplus}$} -- поэлементная сумма.

Обходная связь (\enquote{skip-connection}) ${\mathbf{s}_n^\text{enc}}$ базового блока приводится к числу каналов (\enquote{фичей}) на выходе (при необходимости):
\begin{equation*}
    \mathbf{s}_n^\text{enc} = 
        \begin{cases} 
            \mathbf{W}_n^{\text{enc}} \ast \mathbf{e}_{n-1}, & \text{если} \hspace{0.33em} F_{n-1} \neq F_{n}, \\
            \mathbf{e}_{n-1}, & \text{иначе},
        \end{cases}
\end{equation*}
где ядро проекции ${\mathbf{W}_n^{\text{enc}} \in \mathbb{R}^{F_{n} \times F_{n-1}}}$.

Между слоями энкодера длина тензора уменьшается в 2 раза алгоритмом \enquote{nn.MaxPool1d}:
\begin{equation*}
    \mathbf{e}_{n} \rightarrow \mathbf{e}_{n}{\downarrow_2}.
\end{equation*}

Регуляризация во время обучения выполняется алгоритмом \enquote{nn.Dropout1d} на выходе второй свертки.

Таким образом, в энкодере выполняется последовательное увеличение числа каналов (\enquote{фичей}) с одновременным уменьшением длины тензора.

\subsection{Bottleneck}
\enquote{Bottleneck} реализован на основе трансформера (рис.~\ref{fig:bottleneck_block_diagram}), с предварительным преобразованием размерностей.

\begin{figure}[H]
	\centering
	\includegraphics [width=1.00\linewidth] {parametersEstimation/bottleneck_block_diagram.drawio.png}

	\caption{Блок-схема bottleneck-алгоритма}
	\label{fig:bottleneck_block_diagram}
\end{figure}

На вход \enquote{bottleneck} поступает тензор ${\mathbf{e}_3 \in \mathbb{R}^{F_3 \times L_3}}$ с выхода энкодера. Тензор преобразуется к ${\mathbf{b}_\text{in} \in \mathbb{R}^{F_4 \times L_4}}$:
\begin{equation*}
    \mathbf{b}_\text{in} = \sigma\left(\mathcal{N}^P\left(\mathbf{W}^P \ast \mathbf{e}_3{\downarrow_2}\right)\right),
\end{equation*}
где ${\mathbf{W}^P \in \mathbb{R}^{F_4 \times F_3}}$, и к \enquote{permute}-формату ${\mathbf{b}_\text{in.perm} \in \mathbb{R}^{L_4 \times F_4}}$:
\begin{equation*}
    \mathbf{b}_\text{in} \rightarrow \mathcal{N}^H\left(\mathbf{b}_\text{in}\right) \rightarrow \mathbf{b}_\text{in.perm}
\end{equation*}
для совместимости с популярной реализацией алгоритма \enquote{Multi-Head Self-Attention} (MHSA) \cite{vaswani2017attention}. Тензор разбивается на $h$-\enquote{голов} с количеством каналов ${d_{k}=F_4/h}$ в каждой голове. Результаты вычисления $h$-\enquote{голов}
\begin{equation*}
    \begin{aligned}
        \mathbf{Q}_m &= \mathbf{b}_\text{in.perm} \mathbf{W}_m^Q, \\
        \mathbf{K}_m &= \mathbf{b}_\text{in.perm} \mathbf{W}_m^K, \\
        \mathbf{V}_m &= \mathbf{b}_\text{in.perm} \mathbf{W}_m^V, \\
        \mathbf{H}_m &= \text{Softmax}\left(\frac{\mathbf{Q}_m \mathbf{K}_m^\top}{\sqrt{d_k}}\right) \mathbf{V}_m,
    \end{aligned}
\end{equation*}
где ${m \in \{1, \dots, h\}}$,
\begin{flalign*}
    & \begin{aligned}
        \mathbf{W}_m^Q, \mathbf{W}_m^K, \mathbf{W}_m^V &\in \mathbb{R}^{F_4 \times d_{k}}, \\
        \mathbf{Q}_m, \mathbf{K}_m, \mathbf{V}_m, \mathbf{H}_m &\in \mathbb{R}^{L_4 \times d_k},
    \end{aligned} &
\end{flalign*}
объединяются в единый тензор
\begin{equation*}
    \mathbf{MH}_\text{perm} = \text{Concat}(\mathbf{H}_1, \dots, \mathbf{H}_h)\mathbf{W}^O,
\end{equation*}
где ${\mathbf{W}^O} \in \mathbb{R}^{L_4 \times d_k}$.

${\mathbf{MH}_\text{perm} \in \mathbb{R}^{L_4 \times F_4}}$ получается того же размера, что и ${\mathbf{b}_\text{in.perm}}$.

Результирующий выход алгоритма MHSA формируется в виде суммы со входным тензором:
\begin{equation*}
    \mathbf{b}_\text{attn} = \mathbf{b}_\text{in} \oplus \mathbf{MH},
\end{equation*}
для чего форма представления тензора предварительно преобразуется обратно (${\mathbf{MH}_\text{perm} \rightarrow \mathbf{MH}}$).

После алгоритма MHSA выполняется нормализация данных для их последующей обработки полносвязными слоями (\enquote{Multi-Layer Perceptron}, MLP). Реализованы два \enquote{Feed-Forward}-слоя поканальной обработки с промежуточным увеличением скрытого пространства в $r$-раз:
\begin{equation*}
    \mathbf{z}^\text{MLP} = \mathbf{W}^{2\text{MLP}} \sigma\left(\mathbf{W}^{1\text{MLP}} \mathcal{N}^\text{MLP}\left( \mathbf{b}_\text{attn} \right)\right),
\end{equation*}
\begin{equation*}
    \mathbf{b}_\text{out} = \mathbf{b}_\text{attn} \oplus \mathbf{z}^\text{MLP},
\end{equation*}

\begin{flalign*}
    & \begin{aligned}
        \mathbf{W}^{1\text{MLP}} &\in \mathbb{R}^{(r \cdot F_4) \times F_4}, \\
        \mathbf{W}^{2\text{MLP}} &\in \mathbb{R}^{F_4 \times (r \cdot F_4)},
    \end{aligned} &
\end{flalign*}
где ${r>1}$ -- кратность масштабирования скрытого пространства в MLP-алгоритме.

Регуляризация во время обучения выполняется алгоритмом \enquote{nn.Dropout} после каждого \enquote{Feed-Forward}-слоя.

\subsection{Декодер}
Декодер последовательно приближает данные из скрытого пространства к целевой длине ${L_1}$. В отличие от энкодера, число каналов (\enquote{фичей}) на каждом слое уменьшается с одновременным увеличением длины тензора.

Между слоями декодера длина тензора увеличивается в 2 раза алгоритмом \enquote{nn.Upsample}:
\begin{equation*}
    \mathbf{d}_{n} \rightarrow \mathbf{d}_{n}{\uparrow_2},
\end{equation*}
сопровождающимся последующей сверткой и нормализацией, аналогично входу \enquote{bottleneck}:
\begin{equation*}
    \mathbf{z}_{n}^{\text{dec}} = \sigma\left(\mathcal{N}^{\text{dec}}\left(\mathbf{W}^{\text{dec}} \ast \mathbf{d}_{n}{\uparrow_2}\right)\right),
\end{equation*}
где ${\mathbf{W}^{1\text{dec}}} \in \mathbb{R}^{F_{n-1} \times F_n}$

\enquote{Residual}-архитектура базового блока в точности повторяет блок энкодера (рис.~\ref{fig:residual_conv_diagram}):

\begin{equation*}
    \begin{aligned}
        \mathbf{z}_{n-1}^{1\text{dec}} &= \mathbf{W}_{n-1}^{1\text{dec}} \ast \sigma\left(\mathcal{N}^{2\text{dec}}\left(\mathbf{z}_{n}^{\text{dec}}\right)\right), \\
        \mathbf{z}_{n-1}^{2\text{dec}} &= \mathbf{W}_{n-1}^{2\text{dec}} \ast \sigma\left(\mathcal{N}^{3\text{dec}}\left(\mathbf{z}_{n-1}^{1\text{dec}}\right)\right), \\
        \mathbf{d}_{n-1} &= \mathbf{z}_{n-1}^{2\text{dec}} \oplus \mathbf{s}_{n-1}^\text{dec}.
    \end{aligned}
\end{equation*}
где две свертки с ядрами ${\mathbf{W}_n^{1\text{dec}} \in \mathbb{R}^{2F_n \times F_n}}$, ${\mathbf{W}_n^{2\text{dec}} \in \mathbb{R}^{F_n \times F_n}}$

Слои декодера связаны с соответствующими выходами энкодера обходными связями (\enquote{skip-connections}).


посредством Attention-gated-связей, по аналогии с \cite{oktay2018attention}.

В декодере реализовано глубокое супервизирование (Deep Supervision) \cite{lee2015deeply} с целью более глубокого обучения модели.

\begin{comment}
    \begin{equation*}
    \begin{aligned}
        &\textbf{Encoder:} \\
        \mathbf{e}_1 &= f_{\text{res}}^{(1)}(\mathbf{x}), \quad 
        \mathbf{e}_2 = f_{\text{res}}^{(2)}(\downarrow\mathbf{e}_1), \quad \dots \\
        \\
        &\textbf{Bottleneck:} \\
        \mathbf{b} &= f_{\text{trans}}\left(\sigma(\mathcal{N}(\mathbf{W}_{\text{proj}} \ast \downarrow\mathbf{e}_4))\right) \\
        \\
        &\textbf{Decoder (for } i=4,3,2,1\text{):} \\
        \mathbf{g}_i &= f_{\text{up}}(\mathbf{d}_{i+1}), \quad
        \tilde{\mathbf{e}}_i = f_{\text{att}}(\mathbf{g}_i, \mathbf{e}_i), \quad
        \mathbf{d}_i = f_{\text{res}}^{(i)}(\text{cat}(\mathbf{g}_i, \tilde{\mathbf{e}}_i)) \\
        \\
        &\textbf{Outputs:} \\
        \mathbf{y}_{\text{main}} &= \mathbf{W}_{\text{final}} \ast \mathbf{d}_1 \\
        \mathbf{y}_{\text{ds}}^{(k)} &= \mathbf{W}_{\text{ds}}^{(k)} \ast \mathbf{d}_{k+1} \quad \text{(training only)}
    \end{aligned}
\end{equation*}


\subsection{Attention-Gated Skip Connections}

Допустим, из skip-connection выходит

${\mathbf{x} \in \mathbb{R}^{B \times F_{\text{l}} \times L_{\text{x}}}}$.

Выход декодера

${\mathbf{g} \in \mathbb{R}^{B \times F_{\text{g}} \times L_{\text{g}}}}$.

Выход энкодера сначала согласуется по размерности - приводится к размерности декодера.

${\mathbf{x}' = \text{Interpolate}(\mathbf{x}, \text{size}=L_{\text{g}})}$.

Затем фичи через матрицы весовых коэффициентов ${W_{\text{g}}}$, ${W_{\text{x}}}$ согласуются по числу фичей (приводятся к ${F_\text{int}}$)
\begin{equation*}
    \begin{aligned}
        \mathbf{q} &= \mathcal{N}\left(\mathbf{W}_g \ast \mathbf{g}\right), \quad 
        \mathbf{k} &= \mathcal{N}\left(\mathbf{W}_x \ast \mathbf{x}'\right)
    \end{aligned}
\end{equation*}
где

${\mathbf{q}, \mathbf{k} \in \mathbb{R}^{B \times F_{\text{int}} \times L_{\text{g}}}}$.

\begin{equation*}
    W_g \in \mathbb{R}^{F_{\text{int}} \times F_g \times 1}, \quad
    W_x \in \mathbb{R}^{F_{\text{int}} \times F_l \times 1}, \quad
    W_\psi \in \mathbb{R}^{1 \times F_{\text{int}} \times 1}
\end{equation*}

Суммируются и проходят через функцию активации.

${\mathbf{a} = \sigma_{\text{ReLU}}\left(\mathbf{q} \oplus \mathbf{k}\right)}$.


Преобразуются к одному каналу
${\boldsymbol{\psi} = \sigma_{\text{Sigmoid}}\left(\mathcal{N}\left(\mathbf{W}_{\psi} \ast \mathbf{a}\right)\right)}$

${\boldsymbol{\psi} \in \mathbb{R}^{B \times 1 \times L_{\text{g}}}}$.
    
В декодер поступает 

${\mathbf{y} = \mathbf{x}' \odot \boldsymbol{\psi}}$

${\mathbf{y} \in \mathbb{R}^{B \times F_{\text{l}} \times L_{\text{g}}}}$


${\odot}$ - поэлементное умножение.

${\text{cat}(\cdot,\cdot)}$ - слияние по каналам.

${\downarrow_s}$ - Downsampling by factor $s$ (e.g., MaxPool).

${\uparrow_s}$ - Upsampling by factor $s$ (linear interpolation).

\end{comment}